\documentclass{article}
\usepackage{amsmath, amssymb}
\usepackage{geometry}
\geometry{margin=1in}

\begin{document}

\section*{Riesz Representation Theorem}

\textbf{Important note:} In mathematics, several theorems are known by the name ``Riesz Representation.'' In Cybenko's proof, we are referring to the Riesz Representation Theorem for linear functionals on the space of continuous functions, not merely the well-known Hilbert space version.

\subsection*{Definition of the Space of Continuous Functions}

Let:
\[
K = [0,1]^n
\]
and consider the space:
\[
C(K) = \{ f : K \to \mathbb{R} \mid f \text{ is continuous} \}
\]
We also define the uniform norm:
\[
\| f \|_\infty = \sup_{x \in K} | f(x) |
\]
This normed space is a Banach space.

\subsection*{Formal Statement of the Riesz--Markov--Kakutani Theorem}

Let \( K \) be a compact Hausdorff space, and let:
\[
L : C(K) \to \mathbb{R}
\]
be a continuous linear functional.

Then there exists a unique finite signed regular Borel measure \( \mu \) such that:
\[
L(f) = \int_K f(x) \, d\mu(x)
\]
for every:
\[
f \in C(K)
\]
Moreover:
\[
\| L \| = |\mu|(K)
\]
where \( |\mu| \) is the total variation measure of \( \mu \).

\end{document}